\documentclass[10pt,a4paper]{amsart}
\usepackage[margin=19mm]{geometry}
\usepackage{amsmath,amssymb,mathtools}
\usepackage[scr=rsfs,cal=euler]{mathalfa}
\usepackage{xfrac}
\usepackage[unicode,hidelinks,bookmarks=false]{hyperref}
\newcommand{\ie}{i.e.\ }
\newcommand{\cf}{cf.\ }
\newcommand{\resp}{respectively\ }
\newcommand{\Subsection}{\subsection}
\providecommand{\nobreakdash}{\nobreak\hbox{-}}
\setlength{\emergencystretch}{2em}
\multlinegap0pt
\usepackage{tikz}
\usepackage{amssymb}
\usepackage{mathtools}
\usepackage{enumerate}
\usepackage{xcolor}
\newtheorem{definition}{Definition}
\newtheorem{theorem}{Theorem}
\newtheorem{proposition}{Proposition}
\newtheorem{lemma}{Lemma}
\newcommand{\C}{\mathbb{C}} 
\newcommand{\Domain}[3]{D_{#1, #2, #3}}
\newcommand{\DomainP}[3]{D_{#1, #2, #3}^+}
\newcommand{\Function}[5]{\begin{array}{cccc} #1 : & #2 & \rightarrow & #3 \\ & #4 & \mapsto & #5 \end{array}}
\newcommand{\N}{\mathbb{N}}
\newcommand{\R}{\mathbb{R}}
\newcommand{\T}{\mathbb{T}}
\newcommand{\Z}{\mathbb{Z}} 
\newcommand{\bigo}[1]{\mathcal{O}\left(#1\right)}
\DeclareMathOperator{\id}{id}
\title{Rigidity of Mather's $\beta$-function on a KAM set for analytic billiards-like maps and unique quasianalytic continuation}
\author{Corentin Fierobe, Vadim Kaloshin, Frank Trujillo}
\date{Source: arXiv:2608.15401v1 (CC BY 4.0)}
\begin{document}
\maketitle

\noindent\emph{Source and licence.} Rigidity of Mather's $\beta$-function on a KAM set for analytic billiards-like maps and unique quasianalytic continuation. Version arXiv:2608.15401v1 (CC BY 4.0), distributed by arXiv under the Creative Commons Attribution 4.0 International licence, http://creativecommons.org/licenses/by/4.0/, as recorded in the arXiv licence metadata for this exact version and shown on the abstract page https://arxiv.org/abs/2608.15401v1. Full licence notice: \texttt{licenses/arxiv-2608.15401-CC-BY-4.0.txt}.\par\smallskip

\noindent\textit{Editorial scope.} Counted unit: the article's theorem thm:KAM (source0.tex lines 600--664). The article's complete original proof of it is delivered with it (source0.tex lines 1125--1596, one \texttt{proof} environment of 472 lines, which the article places away from the statement). No mathematics is written, repaired or re-derived: the statement and the proof are the article's own lines, in the article's own notation, and no retained line is substituted or re-flowed.  Retained uncounted as the source-local context the proof needs: lines 260--277; lines 585--588; lines 866--874; lines 1071--1121; lines 1603--1627; lines 2450--2454; lines 2558--2582; lines 2642--2657; lines 2687--2709; lines 2768--2776.  Nothing was omitted from any retained span, so the package carries no omission record.  No retained line contains a citation, so the wrapper carries no bibliography and no bibliography entry.  No retained line contains a figure environment, an image-inclusion command or drawing code, so no image is delivered and the package declares no asset dependency.  Editorial formatting only, in addition: an AMS \texttt{amsart} preamble that restates the article's own declarations and macros used by the retained lines, namely the environments \texttt{definition}, \texttt{theorem}, \texttt{proposition}, \texttt{lemma} on separate counters, together with the standard packages those lines need.  The retained environments print in this document as Definition 1, Theorem 1, Proposition 1, Lemma 1 in order of appearance, whereas the article prints the same items with the numbering of its own full document.  The complete native source of the article as distributed in the arXiv e-print for this exact version is bundled unchanged as \texttt{sources/207742/source0.tex}.
\begin{definition}
\label{def:billiard-like}
A \emph{billiard-like map} is a diffeomorphism onto its image $$T: \T \times (-b, b) \to \T \times \R,$$
for some $b > 0$ and $m \in \N_{\geq 2}$, of the form
\begin{equation}
 \label{equation:expression_billiard_like_maps}
%\Function{T}{\T \times (-b, b)}{\T \times (-b', b')}{(x, y)}{ \big(x+y+\bigo{y^m},y+\bigo{y^{m+1}}\big)},
T(x, y) = \big(x+y+\bigo{y^m},y+\bigo{y^{m+1}}\big),
\end{equation}
 for which there exists a 1-form $\lambda$ of the type
\begin{equation}
    \label{equation:expression_one_form}
\lambda = (y^2 + \mu(x, y))dx + \nu(x, y)dy, \qquad \mu(x, y) = \bigo{y^3}, \quad \nu(x, y)= \bigo{y^2},
\end{equation}
such that $T^*\lambda - \lambda $ is well-defined and exact on $\T \times (-b, b)$. 

The integer $m$ is called the \emph{order} of the billiard-like map $T$.
\end{definition}

\begin{equation}
\label{eq:perturbation_form}
T_R := A+R.
\end{equation}

\begin{equation}
\label{equation:multi_difference_equation}
 \left\{
 \begin{array}{rcl}
 \overline\varphi_1(x+y,y)-\overline\varphi_1(x,y) &=& y(\mathscr T_N \overline R_1(x, y) +\overline\varphi_2(x,y)),\\
 \overline\varphi_2(x+y,y)-\overline\varphi_2(x,y) &=& y (\mathscr T_N \overline R_2(x, y) - [\overline R_2](y)),
 \end{array}
 \right.
\end{equation}

\begin{proposition}
\label{prop:estimates_one_step}
 Let $r,b,h\in(0,1]$, $\sigma \in (0, \tfrac{r}{4}),$ $\overline{\sigma} \in (0, \tfrac{1}{4}\min\{b, h\})$ and $N \in \N_{> 0}$ satisfying
 \begin{equation}
 \label{eq:relation_N_h}
 h \leq \tfrac{1}{2\sqrt 2}\gamma N^{-\tau}.
 \end{equation}
 There exist constants $0 < c_{\ast}< 1 < C_\ast ,$  depending only on $\gamma$ and $\tau$, such that for any map $R\in\mathscr C^1_{r,b,h, 2}(\R^2)$ satisfying
 \begin{equation}
 \label{equation:assumption_R}
 \|R\|_{r,b,h,2}\leq c_\ast
\sigma^{2\tau+3}\overline\sigma^3
 \end{equation}
 for which $T_R$ %= A + R : \Domain{r}{b}{h} \to \C/\Z \times \C$ 
 (as in \eqref{eq:perturbation_form}) is an analytic billard-like map on $\Domain{r}{b}{h}$ (in the sense of Definition \ref{def:billiard-like}) with associated 1-form $\lambda$ defined on $\DomainP{r + \overline \sigma}{b + \overline \sigma}{h + \overline \sigma}$ of the form
\[\lambda(x, y) = (y^2 + y^2\overline \mu(x, y))dx + y \overline \nu(x, y)dy, %\qquad (\mu, \nu) \in \mathscr{C}^{1+}_{r + \sigma, b + \overline \sigma, h + \overline \sigma}, 
\qquad |\overline \mu|_{\DomainP{r + \overline \sigma}{b + \overline \sigma}{h + \overline \sigma}} + |\overline \nu|_{\DomainP{r + \overline \sigma}{b + \overline \sigma}{h + \overline \sigma}} < \frac{1}{4},\]
there exists a diffeomorphism onto its image $\Phi: \DomainP{r - \sigma}{b - \overline \sigma}{h - \overline \sigma} \to \C/\Z \times \C$ such that the following holds:

\begin{enumerate}
\item $R_+ := \Phi^{-1} T_R\Phi - A$ restricted to $\Domain{r-3\sigma}{b-3\overline\sigma}{h-3\overline\sigma}$ %is an element of $\mathscr C^1_{r-3\sigma,b-3\overline\sigma,h-3\overline\sigma, 2}(\R^2)$ satisfying
is well-defined and satisfies
 \[
 \|R_+\|_{r-3\sigma,b-3\overline\sigma,h-3\overline\sigma,2}
 \leq \frac{C_\ast e^{-2\pi N\sigma}}{\pi\sigma\overline \sigma^3} \|R\|_{r,b,h,2} + \frac{C_\ast}{\sigma^{4\tau+5}\overline\sigma^5}\|R\|_{r,b,h,2}^2,
 \]
 % \[
 % \|R_+\|_{r-3\sigma,b-3\overline\sigma,h-3\overline\sigma,2}
 % \leq \frac{e^{-2\pi N\sigma}}{\pi\sigma} \|R\|_{r,b,h,2} + \frac{C_\ast}{\sigma^{4\tau+5}\overline\sigma^5}\|R\|_{r,b,h,2}^2,
 % \]
 \item  $\Phi=I+\varphi$, with $\varphi \in \mathscr C^{1+}_{r-\sigma,b-\overline\sigma,h-\overline\sigma, 1}(\R^2)$ satisfying
 \[
 \|\varphi\|_{r-\sigma,b-\overline\sigma,h-\overline\sigma,1}^+
 \leq 
 \frac{C_2}{\sigma^{2(\tau+1)}\overline\sigma^2}\|R\|_{r,b,h,2},
 \]
 \item $\Phi^{-1}\mid_{\DomainP{r-2\sigma}{b-2\overline\sigma}{h-2\overline\sigma}} = I+\psi$, with  $\psi \in \mathscr C^{1+}_{r-2\sigma,b-2\overline\sigma,h-2\overline\sigma, 1}(\R^2)$ satisfying
 \[
 \|\psi\|^+_{r-2\sigma,b-2\overline\sigma,h-2\overline\sigma,1} 
 \leq 
 \frac{KC_2}{\sigma^{2(\tau+1)}\overline\sigma^2}\|R\|_{r,b,h,2},
 \]
%  \item $\lambda_+ := \Phi^* \lambda \mid_{\DomainP{r - 2\sigma}{b - 2\overline \sigma}{h - 2\overline \sigma}}$ is of the form $\lambda_+(x, y) = (y^2 + \mu_+(x, y))dx + \nu_+(x, y)dy$ with $\mu_+, \nu_+ \in \mathscr{C}^1_{r - 2\sigma, b - 2\overline \sigma, h - 2\overline \sigma}$ satisfying
%  \[ 
%    |\overline \mu_+|_{\DomainP{r - 2 \sigma}{b -  2\overline \sigma}{h - 2\overline \sigma}} + |\overline \nu_+|_{\DomainP{r - 2 \sigma}{b -  2\overline \sigma}{h - 2\overline \sigma}} \leq \left(1 + C\|\varphi\|_{r-\sigma,b-\overline\sigma,h-\overline\sigma,1}^+\right)\left( |\overline \mu|_{\DomainP{r + \overline \sigma}{b + \overline \sigma}{h + \overline \sigma}} + |\overline \nu|_{\DomainP{r + \overline \sigma}{b + \overline \sigma}{h + \overline \sigma}}\right) ,
%  \]
%  for some constant $C > 0$ not depending on any of the parameters,
% \item $S^+ := S \circ \Phi \mid_{\Domain{r-3\sigma}{b-3\overline\sigma}{h-3\overline\sigma}}$ satisfies  $T_{R_+}^* \lambda_+ - \lambda_+ = dS_+.$ 
\end{enumerate}
 where $C_2$ and $K$ are the constants defined at Propositions \ref{proposition:solution_truncated_difference_equation} and \ref{proposition:inverse_estimates} respectively.
\end{proposition}

\begin{proposition}
\label{proposition:solution_truncated_difference_equation}
Let $r,b,h,\sigma,\overline\sigma \in(0,1]$ with $r-\sigma,b-\overline\sigma,h-\overline\sigma>0$, $N \in \N_{> 0}$ and $\overline R_1, \overline R_2\in\mathscr C^1_{r,b,h}(\R)$. 
If 
\[
h \leq \tfrac{1}{2\sqrt 2}\gamma N^{-\tau},
\]
there exist $\overline \varphi_1, \overline \varphi_2\in\mathscr C^{1+}_{r-\sigma,b-\overline\sigma,h-\overline\sigma}(\R)$ satisfying \eqref{equation:multi_difference_equation}. Moreover,
 \begin{equation}
 \label{equation:norm_control_difference_equation}
  \max_{i=1,2}
 \|\overline \varphi_i\|^+_{r - \sigma, b - \overline \sigma, h - \overline \sigma} 
 \leq 
 \frac{C_2}{\sigma^{2(\tau+1)}\overline\sigma^2}
 \max_{i=1,2} 
 \|\overline R_i\|_{r, b, h},
 % \max_{i=1,2}
 % \|\varphi_i\|^+_{\mathscr C_{\text{hol}}^1(K_{h-\overline\sigma}^{b-\overline\sigma},\mathscr C^{\omega}(\T_{r-\sigma},\C))} 
 % \leq 
 % \frac{C_2}{\sigma^{2(\tau+1)}\overline\sigma^2}
 % \max_{i=1,2} 
 % \|R_i\|_{\mathscr C_{\text{hol}}^1(K_h^b,\mathscr C^{\omega}(\T_{r},\C))},
 \end{equation}
where $C_2 = C_2(\gamma, \tau) = \tfrac{\overline{C_2}(\tau)}{\gamma^2}$ and $\overline{C_2}(\tau) > 1$ is a constant depending only on $\tau$.
\end{proposition}

\begin{proposition}
\label{proposition:norm_power_times_map}
Let $b>0$, $K\subset B(0,b)$ be a compact set, $B$ be a Banach space, an integer $s\geq 1$, and $\varphi:K\to B$ be a $\mathscr C^1$-holomorphic map. Then
$$\|y^s\varphi\|_{\mathscr C^1_{\text{hol}}(K,B)}\leq 2sb^{s-1}\|\varphi\|_{\mathscr C^1_{\text{hol}}(K,B)}.$$
\end{proposition}

\begin{proposition}
\label{proposition:composition_estimates}
Let $m\in\Z_{\geq 1}$, $f\in\mathscr C^1_{r,b,h,m}$ and $\varphi=(\varphi_1,\varphi_2):\T_{r-\sigma}\times K_{b-\overline\sigma}^{h-\overline\sigma}\to\C^2$ analytic such that 
\begin{equation}
\label{eq:size_compo}
\max\left\{\|\varphi_1\|_{\mathscr C_{\text{hol}}^1(K_{h-\overline\sigma}^{b-\overline\sigma})}, \|\varphi_2\|_{\mathscr C_{\text{hol}}^1(K_{h-\overline\sigma}^{b-\overline\sigma})} 
\right\}
\leq \kappa\min(\sigma,\overline\sigma)
\end{equation}
where $\kappa,\sigma,\overline\sigma\in(0,1)$ satisfy $r-\sigma,h-\overline\sigma,b-\overline\sigma>0$. 

Then the maps $g,g_x,g_y:\T_{r-\sigma}\times K_{b-\overline\sigma}^{h-\overline\sigma}\to \C^2$ defined by $g=f\circ(I+\varphi)$, $g_x=\partial_xf\circ(I+\varphi)$ and $g_y=\partial_yf\circ(I+\varphi)$ are respectively elements of $\mathscr C^1_{r-\sigma,b-\overline\sigma,h-\overline\sigma,m}$, $\mathscr C^1_{r-\sigma,b-\overline\sigma,h-\overline\sigma,m}$ and $\mathscr C^1_{r-\sigma,b-\overline\sigma,h-\overline\sigma,m-1}$ whose norms satisfy
\begin{equation}
 \label{equation:inequality_composition}
 \|g\|_{r-\sigma,b-\overline\sigma,h-\overline\sigma,m} \leq K_1\|f\|_{r,b,h}
\end{equation}
\[
\|g_x\|_{r-\sigma,b-\overline\sigma,h-\overline\sigma,m}\leq \frac{K_1}{\sigma}\|f\|_{r,b,h}
\]
and
\[
\|g_y\|_{r-\sigma,b-\overline\sigma,h-\overline\sigma,m-1}\leq \frac{K_1}{\overline\sigma}\|f\|_{r,b,h}
\]
where $K_1>0$ is a constant depending uniquely on $m$ and $\kappa$.
\end{proposition}

\begin{proposition}
\label{proposition:difference_estimates}
Let $m,n\in\Z_{\geq0}$, $f\in\mathscr C^1_{r,b,h,m}$ and $\varphi\in\mathscr C^1_{r-\sigma,b-\overline\sigma,h-\overline\sigma,n}$ such that 
\begin{equation}
\label{eq:size_compo_2}
\|\varphi\|_{r-\sigma,b-\overline\sigma,h-\overline\sigma,n} \leq \kappa\min(\sigma,\overline\sigma)
\end{equation}
where $\kappa,\sigma,\overline\sigma\in(0,1)$ satisfy $r-\sigma,h-\overline\sigma,b-\overline\sigma>0$.  

Then the map $h:\T_{r'}\times K_{b'}^{h'}\to \C^2$ defined by
$h=f(I+\varphi)-f$ is an element of $\mathscr C^1_{r-\sigma,b-\overline\sigma,h-\overline\sigma,m+n}$ whose norm satisfies
$$\|h\|_{r-\sigma,b-\overline\sigma,h-\overline\sigma,m+n}\leq K_2\|f\|_{r,b,h,m}\|\varphi\|_{r-\sigma,b-\overline\sigma,h-\overline\sigma,n}\left(
\frac{1}{\sigma}
+\frac{1}{\overline\sigma}\right)$$
where $K_2>0$ is a constant depending only on $m$, $b$ and $\kappa$.
\end{proposition}

\begin{proposition}
\label{proposition:inverse_estimates}
 Let $\Phi$ be a diffeomorphism of the form $\Phi=I+\varphi$ where $\varphi\in\mathscr C^{1+}_{r,b,h,m}$ is a $\mathscr C^1$-holomorphic map. Let $\kappa,\sigma,\overline\sigma\in(0,1)$ satisfy $r-\sigma,h-\overline\sigma,b-\overline\sigma>0$ and assume that 
 \begin{equation}
 \label{equation:assumption_inverse}
 \|\varphi\|_{r,b,h,m}^+
 \leq
 \frac{\kappa}
 {K}\frac{\sigma\overline\sigma}{\sigma+\overline\sigma}
 \end{equation}
 where $K$ is the largest of the constants $K_1$ and $K_2$ given in Propositions \ref{proposition:composition_estimates} and \ref{proposition:difference_estimates}. Then we can write $\Phi^{-1} = I+\psi$ where 
 $\psi\in\mathscr C^{1+}_{r-\sigma,b-\overline\sigma,h-\overline\sigma,m}$
 satifies the estimates
 \begin{equation}
 \label{equation:norm_inverse}
 \|\psi\|_{r-\sigma,b-\overline\sigma,h-\overline\sigma,m}^+\leq K\|\varphi\|_{r,b,h,m}^+
 \end{equation}
 Moreover $\varphi+\psi\in\mathscr C^{1+}_{r-\sigma,b-\overline\sigma,h-\overline\sigma,2m}$ and it holds that
 \begin{equation}
 \label{equation:norm_inverse_plus_diffeo}
 \|\varphi+\psi\|_{r-\sigma,b-\overline\sigma,h-\overline\sigma,2m}^+ \leq K^2({\|\varphi\|^+})^2\left(\frac{1}{\sigma}+\frac{1}{\overline\sigma}\right).
 \end{equation}
\end{proposition}

\begin{lemma}
\label{lem:cauchy}
Let $f \in \mathscr C_{\text{hol}}^1(\overline{K_h},B)$ for some closed set $K \subseteq \C$ without isolated points and some $0 < h < 1$. Then,
\[ \|f^{(r)}\|_{\mathscr C^1_{\text{hol}}(K,B)} \leq \frac{C_r}{h^r} \|f\|_{\mathscr C^1_{\text{hol}}(\overline{K_h},B)},\]
for any $r \geq 1$, where $C_r > 1$ is a constant depending only on $r$. Moreover,
\[
\|f\|_{\mathscr C^1_{\text{hol}}(K,B)}\leq \frac{5}{h}\sup_{\overline{K_h}}|f|.
\]
\end{lemma}

\begin{theorem}
\label{thm:KAM}
There exists a continuous function $\varepsilon: (0,1) \times [1, +\infty) \to (0, 1)$ such that for any $(\gamma,\tau)\in(0,1)\times[1,+\infty)$ the following holds.

Let $r,b\in(0,1)$ and $R: \T_r \times B(0, b) \to \C^2$ be a real analytic map of the form
 \begin{equation}
 \label{eq:form_R}
 R(x, y) = (y^2 \overline R_1(x, y), y^3 \overline R_2(x, y)),
 \end{equation}
 for which $T_R$ (as in \eqref{eq:perturbation_form}) is an analytic billard-like map of order $2$ (in the sense of Definition \ref{def:billiard-like}) satisfying 
 \begin{equation}
 \label{eq:smallness_condition_R}
 \max_{j=1,2}
 \| \overline R_j \|_{\mathscr C^1_{\textup{hol}}(B(0,b))}
 \leq \varepsilon_0(\gamma,\tau,r,b) :=\varepsilon(\gamma,\tau)r^{8\tau+10}
 \min\left\{\left(\frac{r}{\ln r}\right)^{\tau},b\right\}^{10},
 \end{equation}
  with associated 1-form $\lambda$ defined on $\T_{r + 2b} \times B(0, 2b)$ of the form
\begin{equation}
\label{eq:form_form}
\lambda(x, y) = (y^2 + y^2\overline \mu(x, y))dx + y \overline \nu(x, y)dy, 
\end{equation}
satisfying
\begin{equation}
    \label{eq:smallness_condition_form}
    | \overline \mu|_{\T_{r + 2b} \times B(0, 2b)},\, | \overline \nu|_{\T_{r + 2b} \times B(0, 2b)} < \frac{1}{8},
\end{equation}
and 
\begin{equation}
    \label{eq:exact_symplectic}
    T^*\lambda - \lambda = dS,
\end{equation}
for some real-analytic map $S: \T_r \times B(0, b) \to \C$.

Then, there exists a $\mathscr{C}^{\infty}$-holomorphic map
\[
{\varphi}:
\left\{
\begin{array}{ccc}
\mathcal K(\gamma, \tau)
& \to &
\mathscr C^\omega(\T_{r/4}, \C^2)\\
y & \mapsto & (y\overline \varphi_1(\cdot,y),y^2\overline \varphi_2(\cdot,y))
\end{array}\right.
\qquad 
\mathcal K(\gamma, \tau):= K(\gamma, \tau)\cap \overline{B(0,b/2)},
\]
satisfying 
    \begin{equation*}
 \label{eq:bounds_curve_constants}
\max\{
\| \overline \varphi_1\|_{\mathscr C^1_{\textup{hol}}(\mathcal K(\gamma, \tau))},
\|\overline \varphi_2\|_{\mathscr C^1_{\textup{hol}}(\mathcal K(\gamma, \tau))}\}
\leq \tfrac{1}{2},
 \end{equation*}
such that the transformation $\Phi=\id +{\varphi}: \T_{r/2} \times \mathcal K(\gamma, \tau) \to \T_{r/2} \times B(0, b)$ given by
 \[
  \Phi(x,y) := (x+y\overline \varphi_1(x,y),y+y^2\overline \varphi_2(x,y))
 \]
is well-defined, diffeomorphic onto its image, and satisfies
\[
T\Phi(x,y) = \Phi(x+y,y),
\quad \forall (x,y)\in\T_{r/2}\times \mathcal K(\gamma, \tau).
\]
\end{theorem}

\begin{proof}[Proof of Theorem \ref{thm:KAM}]
We elaborate an iterative scheme relying on Proposition \ref{prop:estimates_one_step}. We divide the proof into three steps:
\begin{enumerate}
    \item {\it\underline{Sequences of parameters.}} Given two fixed parameters $\sigma,\overline\sigma\in(0,1)$, we define exponentially decreasing sequences
    $(\sigma_n)_{n\geq 0}$ and $(\overline\sigma_n)_{n\geq 0}$. Given $r,b\in(0,1)$, we can use them to build decreasing sequences
    $(r_n)_{n\geq 0}$, $(b_n)_{n\geq 0}$, $(h_n)_{n\geq 0}$ converging respectively to $r_\infty\geq r/2$, $b_\infty\geq b/2$ and $0$. 
    We construct additional sequences $(\varepsilon_n)_{n\geq 0}$ and $(N_n)_{n\geq 0}$  where $\varepsilon_n$ converges superexponentially to $0$ and $N_n$ is an integer satisfying $h_n\leq \gamma N^{-\tau}/2\sqrt 2$ for any $n$. Moreover $\varepsilon_0$ is of the form $\varepsilon_0 = \varepsilon(\gamma,\tau)r^{8\tau+10}b^{10}$.
    \item {\it\underline{Sequences of maps.}}  We construct sequences of real analytic $\C^2$-valued maps $R_n\in\mathscr C^1_{r_n,b_n,h_n,2}$, $\varphi_n\in\mathscr C^{1+}_{r_n-\sigma_n,b_n-\overline\sigma_n,h_n-\overline\sigma_n, 1},$ $\psi_n\in\mathscr C^{1+}_{r_n-2\sigma_n,b_n-2\overline\sigma_n,h_n-2\overline\sigma_n, 1}$, and of real analytic $\C$-valued maps $\overline \mu_n, \overline \nu_n \in  \mathscr C^{1+}_{r_n + \overline \sigma_n,b_n + \overline \sigma_n,h_n + \overline \sigma_n},$ satisfying the estimates
   % $\|R_n\|\leq \varepsilon_n$, $\|\varphi_n\|\leq \varepsilon_n^{3/4}$, $\|\psi_n\|\leq \varepsilon_n^{3/4}$ 
   \[  \|R_n\|_{r_n, b_n, h_n, 2} \leq \varepsilon_n, \qquad |\overline \mu_n|_{\DomainP{r_n + \overline \sigma_n}{b_n + \overline \sigma_n}{h_n + \overline \sigma_n}} + |\overline \nu_n|_{\DomainP{r_n + \overline \sigma_n}{b_n + \overline \sigma_n}{h_n + \overline \sigma_n}} < \frac{1}{4} \]
   \[\|\varphi_n\|^+_{r_n - \sigma_n, b_n - \overline \sigma_n, h_n - \overline \sigma_n, 1} \leq \varepsilon_n^{3/4},\qquad  \|\psi_n\|^+_{r_n - 2\sigma_n, b_n - 2\overline \sigma_n, h_n - 2\overline \sigma_n, 1}\leq \varepsilon_n^{3/4},\]
   and such that, for any $ n \geq 0$,
 \[ \lambda_{n} = (y^2 + y^2 \overline \mu_{n})dx + y \overline \nu_{n}dy,\]
 satisfies
    \[ 
    T_{R_{n}}^* \lambda_{n} - \lambda_{n} = d S_{n},
    \]
        for some real-analytic function $S_{n}: \Domain{r_n}{b_n}{h_n} \to \C$, and
    \[
    T_{R_{n+1}} = \Phi_n^{-1}\circ T_{R_n}\Phi_n : \Domain{r_{n + 1}}{b_{n + 1}}{h_{n + 1}}  \to \C/\Z \times \C
    \]
    is well-defined, where $\Phi_n = I+\varphi_n$ and $\Phi_n^{-1}\mid_{\DomainP{r_n-2\sigma_n}{b_n-2\overline\sigma_n}{h_n-2\overline\sigma_n}} = I+\psi_n$.    
    
    \item {\it\underline{Conjugating diffeomorphism.}} We show that, for any $k \geq 1$, $\Phi_0\circ\cdots\circ\Phi_n$ converges, as $n$ goes to $+\infty$, to a $\mathscr C^{k}$-holomorphic diffeomorphism $\Phi$ satisfying the conclusions of Theorem \ref{thm:KAM}.
\end{enumerate}

\vspace{0.2cm}
\noindent {\it\underline{Sequences of parameters.}}
We start with $\sigma,\overline\sigma\in(0,1)$ that we are going to set later. Let 
\[
\beta = 4^\tau,
\]
and introduce the sequences $(\sigma_n)_{n\geq 0}$ and $(\overline\sigma_n)_{n\geq 0}$ by setting
\[
\sigma_n = \frac{\sigma}{2^n},
\quad
\overline\sigma_n = \frac{\overline\sigma}{\beta^n},
\qquad\qquad n\in\Z_{\geq 0}.
\]
We first consider the following sequences $(r_n)_{n\geq 0}$, $(b_n)_{n\geq 0}$, $(h_n)_{n\geq 0}$ defined for $n\geq 0$ by
\begin{align*}
r_{n+1} & =  r_n-3\sigma_n, \qquad r_0=r;\\
b_{n+1} & =  b_n-3\overline\sigma_n, \qquad b_0=b;\\
h_{n+1} & =  h_n-3\overline\sigma_n, \qquad h_0=\tfrac{3\beta\overline\sigma}{\beta-1}.
\end{align*}
We easily check the following result:
\begin{lemma}
    \label{lemma:prop_seq_strips}
    The sequence $(h_n)_{n\geq 0}$ is of the form
    \[
    h_n = \frac{h_0}{\beta^n},\qquad n\geq 0
    \]
    hence decreases to $0$.
    Moreover if $\sigma\leq r/6$ and $\overline\sigma\leq \frac{\beta-1}{6\beta}b$, the decreasing sequences $(r_n)_{n\geq 0}$ and $(b_n)_{n\geq 0}$ converge respectively to $r_\infty\geq r/2$ and to $b_\infty\geq b/2$.
\end{lemma}

We now introduce a constant $\tilde c>0$ satisfying simultaneously the following assumptions:
\begin{equation}
    \label{eq:tilde_c}
    \tilde c\leq c_\ast,
    \qquad
    2^{8\tau+10}\beta^{10}\tilde c^{1/2}\leq 1,
    \qquad
    C_\ast(K+1)\tilde c < \tfrac{1}{2},
    \qquad
    \frac{K\tilde c^{3/4}\beta}{\beta-1}\leq\frac{1}{2},
\end{equation}
where $c_\ast$, $C_\ast$, $C_2$ and $K$ are the positive constants appearing in Proposition \ref{prop:estimates_one_step}.
We can define a decreasing sequence $(\varepsilon_n)_{n\geq0}$ of positive numbers by setting
\[
\varepsilon_0 = \tilde c\sigma^{8\tau+10}\overline\sigma^{10}
\quad\text{and}\quad 
\varepsilon_n = \varepsilon_0^{(3/2)^n},
\quad n\geq 0.
\]
Consider then the sequence of positive integers 
$(N_n)_{n\geq 0}$ defined by 

\[
N_n = \left\lceil
\frac{1}{2\pi\sigma_n}
\left|
\ln\left(
\frac{\pi}{2}\sigma_n\varepsilon_n
\right)
\right|
\right\rceil.
\]
We prove simple properties for them:
\begin{lemma}
    \label{lemma:prop_seq_eps_N}
    There exists $\overline\sigma_0\in(0,1)$ of the form
    \[
    \overline\sigma_0 = \sigma_\ast(\gamma,\tau)
    \left(\frac{\sigma}{\ln\sigma}\right)^\tau
    ,\qquad \sigma_\ast(\gamma,\tau)\in(0,1)
    \
    \]
    such that for $\overline\sigma\in(0,\overline\sigma_0]$ the sequences defined above satisfy the following, for any $n\geq 0$:
    \begin{enumerate}
        \item $\varepsilon_n\leq \tilde c \,\sigma_n^{8\tau+10}\overline\sigma_n^{10}$;
        \item $h_n \leq \tfrac{1}{2\sqrt 2}\gamma N_n^{-\tau}$;
        \item $\frac{C_\ast e^{-2\pi N_n\sigma_n}}{\pi\sigma_n\overline\sigma_n^3}\varepsilon_n + \frac{C_\ast}{\sigma_n^{4\tau+5}\overline\sigma_n^5}\varepsilon_n^2\leq \varepsilon_{n+1}$.
        %$\frac{e^{-2\pi N_n\sigma_n}}{\pi\sigma_n}\varepsilon_n+\frac{C_\ast}{\sigma_n^{4\tau+5}\overline\sigma_n^5}\varepsilon_n^2\leq \varepsilon_{n+1}$.
        \item \label{cond:fast_decrease} Moreover, for any $a, k,\ell > 0$, 
        \[
        \lim_{n \to \infty} \frac{\varepsilon_n^a}{\sigma_n^{\ell}\overline\sigma_n^{k}} \to 0.
        \]
    \end{enumerate}
\end{lemma}

\begin{proof}[Proof of Lemma \ref{lemma:prop_seq_eps_N}]
We prove the first assertion by induction, noticing that the case when $n=0$ is immediate by construction of $\varepsilon_0$. Assume that the result holds for a given $n\geq 0$. From the definitions of the sequences $(\varepsilon_n)_{\geq 0}$, $(\sigma_n)_{\geq 0}$ and $(\overline\sigma_n)_{\geq 0}$ follows
\[
\frac{\varepsilon_{n+1}}{\tilde c\sigma_{n+1}^{8\tau+10}\overline\sigma_{n+1}^{10}}
= 
2^{8\tau+10}\beta^{10}\frac{\varepsilon_{n}^{3/2}}{\tilde c\sigma_{n}^{8\tau+10}\overline\sigma_{n}^{10}}
\leq 
2^{8\tau+10}\beta^{10} \frac{\tilde c^{3/2}\sigma_{n}^{12\tau+15}\overline\sigma_{n}^{15}}{\tilde c\sigma_{n}^{8\tau+10}\overline\sigma_{n}^{10}}
\leq 
2^{8\tau+10}\beta^{10}\tilde c^{1/2}
\sigma_{n}^{4\tau+5}\overline\sigma_{n}^{5}\leq 1
\]
where the last inequality follows from the fact that $\sigma_n,\overline\sigma_n\in(0,1)$ and Assumption \eqref{eq:tilde_c} on $\tilde c$. 

To prove the second assertion, we use the definition of $(N_n)_{n\geq 0}$ to observe that for $n\geq 0$,
\begin{align*}
    h_n^{1/\tau}N_n
    & \leq 
    \frac{h_0^{1/\tau}}{\beta^{n/\tau}}
    \left(
    \frac{1}{2\pi\sigma_n}
    \left|
    \ln\left(
    \frac{\pi}{2}\sigma_n \varepsilon_n
    \right)
    \right|
    +1
    \right)\\
    & =
    \frac{h_0^{1/\tau}}{\beta^{n/\tau}}
    \left(
    \frac{2^n}{2\pi\sigma}
    \left|
    \ln\left(
    \frac{\pi\sigma}{2^{n+1}}(\varepsilon_0)^{(3/2)^n}
    \right)
    \right|
    +1
    \right)\\
    & \leq
    \frac{h_0^{1/\tau}}{2\pi\sigma}
    \left(
    \frac{2}{\beta^{1/\tau}}
    \right)^n
    \left|
    \ln\left(
    \frac{\pi\sigma}{2^{n+1}}
    \right)
    \right|
    +
    \frac{h_0^{1/\tau}}{2\pi\sigma}
    \left(
    \frac{3}{\beta^{1/\tau}}
    \right)^n
    \left|
    \ln\left(
    \varepsilon_0
    \right)
    \right|
    +
    \frac{h_0^{1/\tau}}{\beta^{n/\tau}}.
    \\
\end{align*}
By the choice of $\beta$, the first, second and third term converge to $0$ in $n$. 
Moreover the sup of each term can bounded by 
\[
C\overline\sigma^{1/\tau}\frac{\ln(\sigma)}{\sigma}
\]
where $C>0$ is a constant depending only on $\gamma$ and $\tau$.
Hence we can define $\overline\sigma_0\in(0,1)$ as in the statement of Lemma \ref{lemma:prop_seq_eps_N} such that for any $\overline\sigma\in(0,\overline\sigma_0)$ we have
\[
\sup_{n\geq 0} 
    \frac{h_0^{1/\tau}}{2\pi\sigma}
    \left(
    \frac{2}{\beta^{1/\tau}}
    \right)^n
    \left|
    \ln\left(
    \frac{\pi\sigma}{2^{n+1}}
    \right)
    \right|
    +
    \frac{h_0^{1/\tau}}{2\pi\sigma}
    \left(
    \frac{3}{\beta^{1/\tau}}
    \right)^n
    \left|
    \ln\left(
    \varepsilon_0
    \right)
    \right|
    +
    \frac{h_0^{1/\tau}}{\beta^{n/\tau}}
    \leq
    \left(\frac{\gamma}{2\sqrt 2}\right)^{1/\tau}
\]
and the second assertion is proven.

To prove the third assertion, we give an upper bound on the two terms in the sum of the left-hand side of the inequality. First by definition of $N_n$ and \eqref{eq:tilde_c},
\[
    \frac{C_\ast e^{-2\pi N_n\sigma_n}}{\pi\sigma_n \overline \sigma_n^3}\varepsilon_n
    \leq 
    \frac{C_\ast}{\pi\sigma_n\overline \sigma_n^3}
    e^{
    \ln\left(
    \frac{\pi}{2}\sigma_n\varepsilon_n
    \right)
    }
    \varepsilon_n
    =
        \frac{C_\ast}{2 \overline \sigma_n^3}  \varepsilon^{2}_n
        =
    \frac{C_\ast}{2 \overline \sigma_n^3}  \left(\tilde c\sigma_n^{8\tau+10}\overline\sigma_n^{10}\right)^{1/2}\varepsilon^{3/2}_n
    \leq
    \frac{1}{2}\varepsilon_n^{3/2}.
\]
For the second term, we use the first assertion proven above to obtain
\[
    \frac{C_\ast}{\sigma_n^{4\tau+5}\overline\sigma_n^5}\varepsilon_n^2
    \leq
    \frac{C_\ast}{\sigma_n^{4\tau+5}\overline\sigma_n^5}
    \left(\tilde c\sigma_n^{8\tau+10}\overline\sigma_n^{10}\right)^{1/2}
    \varepsilon_n^{3/2}
    = 
    C_\ast\tilde c^{1/2}\varepsilon_n^{3/2}
    \leq 
    \frac{1}{2}\varepsilon_n^{3/2}
\]
where the last inequality comes from the third estimates on $\tilde c$ given in \eqref{eq:tilde_c}. Whence the third assertion is proven. 

Finally the fourth assertion follows using similar arguments.
\end{proof}

\vspace{0.2cm}
\noindent {\it\underline{Sequences of maps.}}
Now we assume that $\sigma =  r/6$ and $\overline\sigma= \min\{\overline\sigma_0,(\beta-1)b/\beta\}
$ so that the Lemmas \ref{lemma:prop_seq_strips} and \ref{lemma:prop_seq_eps_N} are satisfied. We define the desired sequences of maps by induction. Note that with this choice of $\sigma, \overline\sigma$, we can write 
\[
\varepsilon_0 = \varepsilon(\gamma,\tau)r^{8\tau + 10}
\min\left\{\sigma_\ast(\gamma,\tau)\left(\frac{r}{6\ln (r/6)}\right)^{\tau},b\right\}^{10}>0
\]
where $\varepsilon(\gamma,\tau)$ depends only on $\gamma$ and $\tau$. By eventually shrinking $\varepsilon$, one can assume that $\varepsilon_0$ has the form given in \eqref{eq:smallness_condition_R}.

Let $R$, $\lambda$ and $S$ as in Theorem \ref{thm:KAM}. %, namely, $R: \T_{r + 2b} \times B(0, 2b) \to \C$ real analytic of the form \eqref{eq:form_R}, $\lambda$ a 1-form defined on $\T_{r + 2b} \times B(0, 2b)$ of the form \eqref{eq:form_form}, and  $S :\T_r\times B(0,b)\to\C$ real analytic satisfying $T^*\lambda - \lambda = dS.$
We consider the restrictions of $R$ and $S$ to $\Domain{r_0}{h_0}{b_0}$ and of $\lambda$ to $\Domain{r_0 + \overline \sigma_0}{h_0 + \overline \sigma_0}{b_0 + \overline \sigma_0}$,  which we denote by $R_0$, $S_0$ and $\lambda_0 = (y^2 + y^2 \overline \mu_0) dx + y \overline \nu_0 dy$, respectively. We assume that
\[
\|R_0\|_{r_0,b_0,h_0,2}\leq \varepsilon_0, \qquad  |\overline \mu_0|_{\DomainP{r_0 + \overline \sigma_0}{b_0 + \overline \sigma_0}{h_0 + \overline \sigma_0}} + |\overline \nu_0|_{\DomainP{r_0 + \overline \sigma_0}{b_0 + \overline \sigma_0}{h_0 + \overline \sigma_0}} \leq \frac{1}{8},
\]
which follows directly from the bounds of the form \eqref{eq:smallness_condition_R} and \eqref{eq:smallness_condition_form} as in the statement of Theorem \ref{thm:KAM}. Notice that, with these assumptions, we are in position to apply Proposition \ref{prop:estimates_one_step}.

 We will construct the sequences $R_n$, $\lambda_n$ and $S_n$ inductively using Proposition \ref{prop:estimates_one_step}, where $\Phi_{n}$ will be given by the proposition and we define 
 \[R_{n + 1} = \Phi_n^{-1} T_{R_n} \Phi_n - A, \qquad \lambda_{n + 1} = \Phi_n^{*} \lambda_n, \qquad S_{n + 1} = S_n \circ \Phi_n,\]
 restricted to appropriate domains. 
 
 Let $n\geq 0$ and assume that $R_n \in \mathscr C^1_{r_n,b_n,h_n, 2}(\R^2)$, $\lambda_n = (y^2 + y^2 \overline \mu_n) dx + y \overline \nu_n dy$, with $\overline \mu_n, \overline \nu_n \in  \mathscr C^{1+}_{r_n + \overline \sigma_n,b_n + \overline \sigma_n,h_n + \overline \sigma_n}(\R)$, and $S_n \in \mathscr C^1_{r_n,b_n,h_n}(\R)$ have been defined (by iterating Proposition \ref{prop:estimates_one_step}) and satisfy
\begin{equation}
    \label{eq:inductiion_hyp}
\|R_n\|_{r_n,b_n,h_n,2}\leq \varepsilon_n,  \qquad T^* \lambda_n - \lambda_n = dS_n, \qquad \Lambda_n  < \frac{1}{4},
\end{equation}
where, for the sake of clarity, we denote
\begin{equation}
\label{eq:Lambda}
\Lambda_n:=  |\overline \mu_n|_{\DomainP{r_n + \overline \sigma_n}{b_n + \overline \sigma_n}{h_n + \overline \sigma_n}} + |\overline \nu_n|_{\DomainP{r_n + \overline \sigma_n}{b_n + \overline \sigma_n}{h_n + \overline \sigma_n}}. % \qquad \text{ for any } n\geq 0.
\end{equation}

Let us show that, with these assumptions, we can apply Proposition \ref{prop:estimates_one_step} to define $R_{n + 1}$, $\lambda_{n + 1}$ and $S_{n + 1}$  satisfying \eqref{eq:inductiion_hyp} (replacing $n$ by $n + 1$), and thus that we can iterate this process indefinitely.

Indeed by Lemma \ref{lemma:prop_seq_eps_N} and the first assumption on $\tilde c$ given in \eqref{eq:tilde_c}, 
\[
\varepsilon_n\leq \tilde c\sigma_n^{8\tau+10}\overline\sigma_n^{10}\leq 
c_\ast\sigma_n^{2\tau+3}\overline\sigma_n^{3},
\]
and the second assertion of Lemma \ref{lemma:prop_seq_eps_N} states that
\[
h_n\leq\frac{1}{2\sqrt 2}\gamma N_n^{-\tau}.
\]

Hence, Proposition \ref{prop:estimates_one_step} implies the existence of two $\C^2$-valued maps $\varphi_n\in\mathscr C^{1+}_{r_n-\sigma_n,b_n-\overline\sigma_n,h_n-\overline\sigma_n}$ and $\psi_n\in\mathscr C^{1+}_{r_n-2\sigma_n,b_n-2\overline\sigma_n,h_n-2\overline\sigma_n}$ such that
$\Phi_n := I+\varphi_n$ is a diffeomorphism onto its image, $\Phi^{-1}_n\mid_{\DomainP{r_n-2\sigma_n}{b_n-2\overline\sigma_n}{h_n-2\overline\sigma_n}} = I + \psi_n$, and the map 
\[
\Phi^{-1}_nT_{R_n}\Phi_n :\Domain{r_{n + 1}}{b_{n + 1}}{h_{n + 1}}  \to\C/\Z\times \C
\]
is a well-defined real analytic map which can be written as
$\Phi^{-1}_nT_{R_n}\Phi_n = A+R_{n+1}$ for some $R_{n+1}\in\mathscr C^1_{r_{n+1},b_{n+1},h_{n+1}}$.
Moreover, we have the following estimates:
\begin{align*}
\|R_{n+1}\|_{r_{n+1},b_{n+1},h_{n+1},2} 
&\leq
\frac{C_\ast e^{-2\pi N_n\sigma_n}}{\pi\sigma_n\overline \sigma_n^3}\|R_n\|_{r_n,b_n,h_n,2}+\frac{C_\ast}{\sigma_n^{4\tau+5}\overline\sigma_n^5}\|R_n\|_{r_n,b_n,h_n,2}^2\\
&\leq
\frac{C_\ast e^{-2\pi N_n\sigma_n}}{\pi\sigma_n\overline \sigma_n^3}\varepsilon_n+\frac{C_\ast}{\sigma_n^{4\tau+5}\overline\sigma_n^5}\varepsilon_n^2\\
&\leq \varepsilon_{n+1},
\end{align*}
where from the second to third line we used the third assertion of Lemma \ref{lemma:prop_seq_eps_N}. 

By Proposition \ref{prop:estimates_one_step} and the first assertion of Lemma \ref{lemma:prop_seq_eps_N}, we also have the following estimates for the norms of $\varphi_n$ and $\psi_n$:
\begin{align*}
\|\varphi_{n}\|^+_{r_{n}-\sigma_n,b_{n}-\overline\sigma_n,h_{n}-\overline\sigma_n,1} 
&\leq
\frac{C_\ast}{\sigma_n^{2(\tau+1)}\overline\sigma_n^2}\varepsilon_n\\
&\leq
\frac{C_\ast}{\sigma_n^{2(\tau+1)}\overline\sigma_n^2}(\tilde c\sigma_n^{8\tau+10}\overline\sigma_n^{10})^{1/4}
\varepsilon_n^{3/4}\\
&\leq
C_\ast\tilde c^{3/4}
\varepsilon_n^{3/4}\\
&\leq \varepsilon_{n}^{3/4},
\end{align*}
where the last inequality follows from the second assumption on $\tilde c$ given in \eqref{eq:tilde_c}. 
In a similar way we obtain
\[
\|\psi_{n}\|^+_{r_{n}-2\sigma_n,b_{n}-2\overline\sigma_n,h_{n}-2\overline\sigma_n,1} \leq \varepsilon_n^{3/4}.
\]

Noticing that $\varepsilon_n^{3/4} < \sigma_n \overline \sigma_n$, it follows that the maps $\Phi_n$, and $T_{R_{n + 1}}$ are well-defined transformations with the following domains/codomains
\[ 
\Phi_n: \DomainP{r_n - \sigma_n}{b_n - \overline \sigma_n}{h_n - \overline \sigma_n} \to  \DomainP{r_n}{b_n}{h_n},\]
\[T_{R_{n + 1}}: \Domain{r_{n + 1}}{b_{n + 1}}{h_{n + 1}} \to \Domain{r_{n + 1} + \overline \sigma_{n + 1}}{b_{n + 1} + \overline \sigma_{n + 1}}{h_{n + 1} + \overline \sigma_{n + 1}}.
%\Phi_n^{-1}: \T_{r_{n} - 2\sigma_n}\times K_{h_{n} - 2\overline \sigma_n}^{b_{n} - 2\overline \sigma_n} \to \T_{r_{n} - \sigma_n}\times K_{h_{n} - \overline \sigma_n}^{b_{n} - \overline \sigma_n},
\]
In particular, 
\[\lambda_{n + 1}= \Phi_n^* \lambda_n,\qquad S_{n + 1}= S_n \circ \Phi_n,\] are well-defined on $\DomainP{r_{n + 1} + \overline \sigma_{n + 1}}{b_{n + 1} + \overline \sigma_{n + 1}}{h_{n + 1} + \overline \sigma_{n + 1}}$ and $\Domain{r_{n + 1}}{b_{n + 1}}{h_{n + 1}}$, respectively, and satify
\[ T_{R_{n + 1}}^* \lambda_{n + 1} - \lambda_{n + 1} = dS_{n + 1}, \]
on $\Domain{r_{n + 1}}{b_{n + 1}}{h_{n + 1}}$. Denoting 
\[\varphi_n =  (\varphi_{n , 1}, \varphi_{n, 2}) = (y \overline \varphi_{n , 1}, y^2 \overline \varphi_{n, 2}),\]
a direct calculation yields
\begin{align*}
    \lambda_{n + 1} & = y^2\left( (1 + y \overline \varphi_{n, 2})(1 + \overline \mu_n \circ \Phi_n)(1 + \partial_x \varphi_{n, 1}) + y(1 + y \overline \varphi_{n, 2})\overline (\nu_n \circ \Phi_n) \partial_x \overline \varphi_{n, 2}\right)dx \\
    & \quad +y\left(y(1 + y \overline \varphi_{n, 2})(1 + \overline \mu_n \circ \Phi_n)\partial_y \varphi_{n, 1} + (1 + y \overline \varphi_{n, 2})\overline (\nu_n \circ \Phi_n) (1 + \partial_y \varphi_{n, 2}\right)dy,
\end{align*}
Thus, we can express $\lambda_{n + 1}$ as 
\[ \lambda_{n + 1} = (y^2 + y^2 \overline \mu_{n + 1}) dx + y \overline \nu_{n + 1} dy\]
with $\overline \nu_{n + 1}, \overline \mu_{n + 1} \in  \mathscr C^{1+}_{r_{n + 1} + \overline \sigma_{n + 1},b_{n + 1} + \overline \sigma_{n + 1},h_{n + 1} + \overline \sigma_{n + 1}}(\R)$, and there exists a constant $C > 0$, independent of all parameters, such that
\[ \Lambda_{n + 1} \leq \Lambda_n\left(1 + \frac{C}{\overline \sigma_n}\|\varphi_{n}\|^+_{r_{n}-\sigma_n,b_{n}-\overline\sigma_n,h_{n}-\overline\sigma_n,1}\right),  \]
where $\Lambda_n$, $\Lambda_{n + 1}$ are given by \eqref{eq:Lambda}. 

Hence, using the estimates above for $\|\varphi_{n}\|_{r_{n}-\sigma_n,b_{n}-\overline\sigma_n,h_{n}-\overline\sigma_n,1}$, we have
\[ \Lambda_{n + 1} \leq \Lambda_0 \exp\left(C \sum_{j = 0}^n \varepsilon_0^{\tfrac{1}{2}\big(\tfrac{3}{2}\big)^j}\right).\]
Notice that, up to choosing the constant $\tilde c$ in the definition of $\varepsilon$ smaller, we may assume that the sum in the equation above is smaller than $\frac{1}{2C}$. Thus, since $\Lambda_0 < \tfrac{1}{8}$ by assumption, we have
\[ \Lambda_{n + 1} < \Lambda_0  \exp\big(\tfrac{1}{2}\big)< \frac{1}{4}.\]

Therefore, we have shown that $R_{n + 1}$, $\lambda_{n + 1}$ and $S_{n + 1}$ (defined using Proposition \ref{prop:estimates_one_step}) satisfy \eqref{eq:inductiion_hyp} (replacing $n$ by $n + 1$) and thus we can define the sequences described above for all $n \geq 0$.

\vspace{0.2cm}
\noindent {\it\underline{Conjugating diffeomorphism.}}
Consider the composition
\[
\Phi^{(n)} = \Phi_0\circ\Phi_1\circ\cdots\circ\Phi_n,
\qquad n\in\Z_{\geq 0}
\]
and set 
\[
K_{\infty} = \bigcap_{n\geq 0} K_{h_n}^{b_n} = K_0^{b_\infty}.
\]
The following lemma shows that $\Phi^{(n)}:\Domain{r_n-\sigma_n}{h_n-\overline\sigma_n}{b_n-\overline\sigma_n}\to\C/\Z\times\C$ is well-defined and satisfies crucial properties to finish the proof:
\begin{lemma}
    \label{lemma:composition_diffeos}
    For any $n\geq 0$, $\Phi^{(n)}:\Domain{r_n-\sigma_n}{h_n-\overline\sigma_n}{b_n-\overline\sigma_n} \to\C/\Z\times \C$ is a well-defined real-analytic diffeomorphism onto its image which  can be written as
    \[
    \Phi^{(n)} = I + \varphi^{(n)},
    \qquad \varphi^{(n)} \in\mathscr C^1_{r_n-\sigma_n, b_n-\overline\sigma_n, h_n-\overline\sigma_n}.
    \]
    It satisfies
    \begin{enumerate}
        \item $\sup_{n\geq 0} \|\varphi^{(n)}\|_{r_n-\sigma_n, b_n-\overline\sigma_n, h_n-\overline\sigma_n,1}\leq \frac{1}{2}$;
        \item For any integers $n\geq 1$ and $k\geq 0$, there is a constant $C_k>0$ depending only on $k$ such that
        \[
        \|D^k\varphi^{(n)}-D^k\varphi^{(n-1)}\|_{\mathscr C^1_{\text{{hol}}}(K_{\infty})}
        \leq
        C_k\frac{\varepsilon_n^{3/4}}{\sigma_n^{r+1}\overline\sigma_n}.
        \]
    \end{enumerate}
\end{lemma}

\begin{proof}[Proof of Lemma \ref{lemma:composition_diffeos}]
    Given $n\geq 0$, we can formaly write $\varphi^{(n)}$ as
    \[
    \varphi^{(n)} = \sum_{j=0}^n u_j,
    \qquad
    u_j := \varphi_j\circ(I+\varphi_{j+1})\circ\cdots\circ (I+\varphi_n).
    \]
    Let us prove by induction on $j$ dreasing from $n$ to $0$ that
    \begin{equation}
    \label{eq:induction_uj}
    \sum_{i=j+1}^n \|u_j\|_{r_n-\sigma_n, b_n-\overline\sigma_n, h_n-\overline\sigma_n,1}\leq \frac{1}{2}\overline\sigma_j,
    \end{equation}
    and the first assertion will follow by taking $j=0$. The result is true at $j=n$ since by construction $u_n=\varphi_n$, and by Lemma \ref{lemma:prop_seq_eps_N}
    \[
    \|\varphi_n\|_{r_n-\sigma_n, b_n-\overline\sigma_n, h_n-\overline\sigma_n,1}
    \leq \varepsilon_n^{3/4} 
    \leq \tilde c
    \,\sigma_n^{8\tau+10}\overline\sigma_n^{10}
    \leq \frac{1}{2}\overline\sigma_n,
    \]
    where the last inequality follows from the choice of $\tilde c$ and the fact that $\sigma,\overline\sigma\in(0,1)$.
    Assume now that \eqref{eq:induction_uj} is satisfied for a given integer $j$ between $0$ and $n-1$. 
    We can apply Proposition \ref{proposition:composition_estimates} to each $u_i$ with $i\geq j$ and $\kappa=1/2$, noticing that
    \[
    u_i = \varphi_i\circ\left(I+\sum_{i'=i+1}^nu_{i'}\right).
    \]
    We get
    \[
    \|u_i\|_{r_n-\sigma_n, b_n-\overline\sigma_n, h_n-\overline\sigma_n,1}
    \leq K\|\varphi_i\|_{r_n-\sigma_n, b_n-\overline\sigma_n, h_n-\overline\sigma_n,1}
    \leq 
    K\varepsilon_i^{3/4}
    \leq K\tilde c^{3/4}\overline\sigma_i \leq \frac{K\tilde c^{3/4}}{\beta^i}.
    \]
    It implies that
    \[
    \sum_{i=j}^n \|u_i\|_{r_n-\sigma_n, b_n-\overline\sigma_n, h_n-\overline\sigma_n,1}
    \leq 
    K\tilde c^{3/4}\sum_{i=j}^n\frac{1}{\beta^i}
    \leq 
    \frac{K\tilde c^{3/4}\beta}{\beta-1}\overline\sigma_j\leq \frac{1}{2} \overline\sigma_j,
    \]
    where the last inequality follows by the assumptions \eqref{eq:tilde_c} on $\tilde c$. The proof of the first assertion is immediate using this result.

    To prove the second assertion, we apply Lemma \ref{lem:cauchy} so that it suffices to show the statement with $r=0$. Indeed, for $n\geq 1$ we can write
    \[
    \varphi^{(n)}-\varphi^{(n-1)}
    =
    \varphi_{n}
    +
    \varphi^{(n-1)}\circ(I+\varphi_{n})-\varphi^{(n-1)}.
    \]
    Applying Proposition \ref{proposition:difference_estimates}, we obtain
    \[
    \|\varphi^{(n)}-\varphi^{(n-1)}\|
    \leq
    \|\varphi_n\|
    +
    K\|\varphi^{(n-1)}\|\|\varphi_n\|
    \left(\frac{1}{\sigma_{n-1}}+\frac{1}{\overline\sigma_{n-1}}\right)
    \leq 
    \varepsilon_n^{3/4}+K\frac{\varepsilon_n^{3/4}}{\sigma_{n-1}\overline\sigma_{n-1}},
    \]
    where 
    $\|\cdot\|$ stands for $\|\cdot\|_{r_n-\sigma_n, b_n-\overline\sigma_n, h_n-\overline\sigma_n,1}$ to simplify. As a result
    \[
    \|\varphi^{(n)}-\varphi^{(n-1)}\|_{\mathscr C^1_{\text{hol}}(K_\infty)} \leq 2(K+1)\frac{\varepsilon_n^{3/4}}{\sigma_{n-1}\overline\sigma_{n-1}},
    \]
    where we have used Proposition \ref{proposition:norm_power_times_map} and the fact that $\sigma_{n-1},\overline\sigma_{n-1}\in(0,1)$.
\end{proof}

It follows from Lemma \ref{lemma:composition_diffeos} that the sequence of maps $\Phi^{(n)}$ converge to a $\mathscr C^{\infty}$-holomorphic map
\[
\Phi:\T_{r_{\infty}}\times K_\infty\to\C/\Z\times \C,
\]
where the notion of convergence is in the $\mathscr C^k$-holomorphic norm on $K_{\infty}$ for any $k\geq0$. Moreover $\Phi$ is a diffeomorphism onto its image satisfying
\[
\|\Phi-I\|_{\mathscr C^1_{\text{hol}}(K_\infty)}\leq \frac{1}{2}
\]
and
\[
T_R\circ\Phi(x,y)=\Phi\circ A(x,y),
\qquad
(x,y)\in\T_{r_\infty}\times K_\infty.
\]
This concludes the proof of Theorem \ref{thm:KAM}.
\end{proof}

\end{document}
